\documentclass{article}

\usepackage{microtype}
\usepackage{graphicx}
\usepackage{subfigure}
\usepackage{booktabs}
\usepackage{hyperref}
\newcommand{\theHalgorithm}{\arabic{algorithm}}
\usepackage[accepted]{icml2025}
\usepackage{amsmath}
\usepackage{amssymb}
\usepackage[capitalize,noabbrev]{cleveref}
\usepackage[textsize=tiny]{todonotes}

\icmltitlerunning{Predicting Bus Stop Passenger Volume in Singapore}

\begin{document}

\twocolumn[
  \icmltitle{Predicting Bus Stop Passenger Volume in Singapore\\
    from Spatial and Network Features}

  \begin{icmlauthorlist}
    \icmlauthor{Joshua Gloor}{nus}
  \end{icmlauthorlist}

  \icmlaffiliation{nus}{ST5227 Applied Statistical Learning, National University of Singapore}

  \icmlcorrespondingauthor{Joshua Gloor}{joshua@gloor.fyi}

  \icmlkeywords{regression, public transit, feature engineering, spatial data}

  \vskip 0.3in
]

% \printAffiliationsAndNotice{\icmlEqualContribution}  % commented out; ICML cannot
\footnotetext[1]{ST5227 Applied Statistical Learning, National University of Singapore.}

\begin{abstract}
  We predict passenger volume at Singapore bus stops from spatial and network
  features derived from public datasets on residential density, MRT
  connectivity, points of interest, and bus-line structure. Using engineered
  per-stop features, we compare linear and tree-based models
  across four temporal targets (weekday morning, midday, evening, and
  weekend/holiday averages) under repeated train/test splits.
\end{abstract}

\section{Introduction}
\label{sec:intro}

Public bus networks carry the majority of public-transport trips in Singapore,
and understanding where ridership concentrates is central to operational and
planning decisions: where to add capacity, where to open new interchanges,
and where to consolidate under-used stops. Passenger counts are expensive to
collect directly, but a large amount of information about a stop's
surroundings is freely available --- residential density, proximity to rail,
nearby amenities, and the stop's role in the bus network itself. It is
therefore useful to know how well such features, on their own, predict the
volume a stop will see.

We treat this as a supervised regression problem. For each stop we construct
thirteen features from five public datasets and predict four summary volumes
per stop: mean hourly boardings-plus-alightings on weekday mornings,
weekday middays, weekday evenings, and weekend/holiday daytimes. The four
targets let us ask a secondary question: do the same features drive
ridership in different parts of the week, or does the relationship change
between commute and leisure periods?

The project has three goals, aligned with the questions that a transit
operator would naturally ask: (i) \emph{predictive}, which class of models
fits this problem best; (ii) \emph{inferential}, which covariates carry the
signal and whether the data-generating mechanism looks linear; and (iii)
\emph{practical}, whether a simpler, interpretable model is sufficient or
whether the accuracy gap to a black-box model is worth the loss of
interpretability.

\section{Data}
\label{sec:data}

\paragraph{Sources.}
We use the LTSG dataset \citep{rs14153579}, which bundles five tables
relevant to Singapore transit: hourly bus passenger counts
(\texttt{bus\_vol}), bus-line routing (\texttt{bus\_line}), residential
block attributes (\texttt{hdb}), MRT station locations (\texttt{mrt}), and
geocoded points of interest (\texttt{poi}). LTSG does not include
coordinates for bus stops, so we join them against the LTA Bus Stop dataset
on data.gov.sg \citep{lta2025busstop}. Of $\sim$5,000 stops in \texttt{bus\_vol}, 23 ($<0.5\%$)
have no match in the current LTA data and are dropped; this is plausibly
due to stops being retired between the two datasets' collection periods.
After the join and requiring complete targets (see below), $n = 4985$ stops remain.

\paragraph{Target construction.}
The raw \texttt{bus\_vol} table gives hourly tap-in (\texttt{in}) and
tap-out (\texttt{out}) counts per stop, labelled weekday (WD) or holiday/
weekend (H). Visualisation of the hourly distribution
(\cref{fig:hourly-profile}) shows two clear phenomena: a right-skewed
distribution across stops at every hour, and sharply different temporal
patterns on weekdays versus weekends. We therefore define volume as
$\text{in}+\text{out}$ and construct four targets per stop by averaging
volume over four hour ranges:

\begin{itemize}
  \item \texttt{wd\_am}: weekday, 06:00--08:59 (morning peak);
  \item \texttt{wd\_pm}: weekday, 16:00--19:59 (evening peak);
  \item \texttt{wd\_midday}: weekday, 09:00--15:59 (off-peak);
  \item \texttt{h\_avg}: weekend/holiday, 06:00--22:59 (daytime).
\end{itemize}

The evening peak is extended to 19:59 to reflect later typical working
hours locally. All four are continuous scalars, so the supervised task is
regression.

\paragraph{Feature construction.}
We engineer $p = 13$ features per stop, all computable from the stop's
coordinates and the five source tables. Spatial features use a Haversine
distance and are computed against stop centroids. \cref{tab:features}
summarises the feature set; rationale for each is given in the EDA
notebook. Radii (500\,m for HDB and POI, 1\,km for MRT) reflect typical
walking distances in Singapore: longer for rail because riders walk further
to stations.

\begin{table}[t]
  \centering
  \small
  \caption{Per-stop features. All counts and sums are computed within the
    stated radius of the stop. \texttt{is\_terminal} is binary; all other
    features are non-negative reals or counts.}
  \label{tab:features}
  \begin{tabular}{@{}lll@{}}
    \toprule
    Name                           & Source    & Description                  \\
    \midrule
    \texttt{hdb\_units\_500m}      & HDB       & dwelling units               \\
    \texttt{hdb\_rentals\_500m}    & HDB       & rental units                 \\
    \texttt{hdb\_commercial\_500m} & HDB       & commercial blocks            \\
    \texttt{dist\_nearest\_mrt}    & MRT       & metres to nearest station    \\
    \texttt{mrt\_within\_1km}      & MRT       & MRT stations within 1\,km    \\
    \texttt{poi\_education\_500m}  & POI       & schools, universities        \\
    \texttt{poi\_food\_500m}       & POI       & restaurants, cafes, bars     \\
    \texttt{poi\_health\_500m}     & POI       & clinics, hospitals           \\
    \texttt{poi\_leisure\_500m}    & POI       & parks, museums, attractions  \\
    \texttt{poi\_shopping\_500m}   & POI       & retail, malls, supermarkets  \\
    \texttt{poi\_worship\_500m}    & POI       & religious sites              \\
    \texttt{n\_lines}              & bus\_line & distinct bus lines at stop   \\
    \texttt{is\_terminal}          & bus\_line & first/last stop on any route \\
    \bottomrule
  \end{tabular}
\end{table}

\paragraph{Missingness.}
There are no null values in the raw feature tables. However,
\texttt{bus\_vol} does not record zero-volume hours explicitly --- a stop
with no activity during a given hour simply has no row for that cell.
A small number of stops are missing observations for at least one
target bucket --- 15 stops are missing weekday midday data, 13 are
missing weekday morning, and a handful have gaps in the other buckets.
These partial rows are dropped, leaving $n = 4985$ stops with complete
targets for modelling.

\paragraph{Subsetting.}
No further subsetting is applied. The dataset size (a few thousand stops
with thirteen features) is well within what all methods in the course
handle without downsampling.

\begin{figure}[t]
  \centering
  \includegraphics[width=0.95\columnwidth]{../reports/step_plot}
  \caption{Mean hourly tap-in (blue) and tap-out (orange) across stops, with
    shaded IQR bands, on weekdays (left) and weekends/holidays (right).
    Weekday peaks are pronounced and directionally asymmetric; weekend
    volume is flatter and smaller. Motivates the four-target split.}
  \label{fig:hourly-profile}
\end{figure}

\section{Methods}
\label{sec:methods}

We compare a set of regression models chosen to span the main mechanisms
the course covered: pure linear, regularised linear, and two tree ensembles.
Running a range from rigidly linear to
flexibly nonlinear lets us see directly how much structure the data
supports.

\paragraph{Models.}
\begin{itemize}
  \item \textbf{Linear regression} (OLS). Baseline; gives a reference for
        how much of the variance is explained by simple additive effects.
  \item \textbf{Ridge regression} with $\ell_2$ penalty. Guards against
        ill-conditioning from correlated features (e.g.\ HDB units and rentals) and provides a fair linear benchmark.
  \item \textbf{Random forest}. Flexible nonlinear, captures
        interactions. Standard black-box reference.
  \item \textbf{Gradient boosted trees}. Typically the strongest tabular
        baseline; tests whether an ensemble boosting approach materially
        beats a random forest here.
\end{itemize}

Each model is fit independently on each of the four targets; we do not
attempt joint multi-output training (this is discussed as a limitation in
\cref{sec:limitations}).

\paragraph{Evaluation protocol.}
We report 5-fold cross-validated performance, which gives a distribution
of five test scores per model per target rather than a single number.
Folds are shuffled once with a fixed seed, and the same fold indices are
used for every model, so pairwise comparisons are on matched test sets.
The primary metric is test RMSE on the natural passenger-count scale;
mean absolute error (MAE) is reported alongside as a robustness check.
Because targets are right-skewed, models are trained on $\log(1+y)$ and
predictions are inverse-transformed before scoring.

Statistical significance of pairwise model differences is assessed with
a paired $t$-test across folds: paired because models are evaluated on
identical folds, which removes fold-level variance (some folds contain
harder test stops than others). A small $p$-value indicates consistent
outperformance across folds rather than luck on particular splits.

\paragraph{Hyperparameters.}
Ridge $\alpha$ is selected by \texttt{RidgeCV} over 20 values logspaced in
$[10^{-2}, 10^{3}]$, using its built-in leave-one-out criterion on the
outer training fold. For the tree ensembles, a nested 3-fold
cross-validation with \texttt{GridSearchCV} selects random forest's
\texttt{min\_samples\_leaf} from $\{5,10,20\}$, and gradient boosting's
\texttt{learning\_rate} from $\{0.05,0.1\}$ and \texttt{max\_depth} from
$\{\text{None},5,8\}$. Random-forest tree count is fixed at 200, which
is well into the diminishing-returns regime for this dataset size;
gradient boosting uses the library's default of 100 trees with early
stopping, so its effective count is data-dependent.

\section{Results}
\label{sec:results}

\paragraph{Gradient boosting wins on every target.}
\cref{tab:main-results} reports test RMSE across all four models and
targets. Gradient boosting is the lowest-error model on every target,
with random forest a close second, ridge third, and linear regression
fourth. The ranking is identical across targets. Relative to the linear
baseline, gradient boosting reduces RMSE by $36\%$ on
\texttt{wd\_am}, $35\%$ on \texttt{wd\_pm}, $54\%$ on
\texttt{wd\_midday}, and $47\%$ on \texttt{h\_avg}.

\begin{table}[t]
  \centering
  \small
  \caption{Test RMSE (mean over 5 folds) by model and target; lower is
    better. Bold: best per target.}
  \label{tab:main-results}
  \begin{tabular}{@{}lcccc@{}}
    \toprule
    Model          & \texttt{wd\_am}   & \texttt{wd\_pm}   & \texttt{wd\_midday} & \texttt{h\_avg}  \\
    \midrule
    Linear         & 4804.4            & 4667.0            & 4357.5              & 1466.8           \\
    Ridge          & 4679.9            & 4550.3            & 4201.7              & 1416.0           \\
    Random forest  & 3232.8            & 3299.4            & 2166.9              & 868.1            \\
    Gradient boost & $\mathbf{3079.0}$ & $\mathbf{3052.3}$ & $\mathbf{2019.0}$   & $\mathbf{771.3}$ \\
    \bottomrule
  \end{tabular}
\end{table}

\paragraph{The mechanism is strongly nonlinear.}
The gap between linear and tree-based models is the dominant effect in
the comparison: $\sim$1500--2300 RMSE points across all four targets,
corresponding to a roughly $35$--$55\%$ error reduction. That the tree
models also interact features (gradient boosting and random forest
both split on thresholds and combine splits multiplicatively) is
consistent with the mechanism being multivariate-nonlinear, not a
failure to include the right linear terms.

Interestingly, the paired $t$-test does \emph{not} declare the
tree-vs-linear gaps significant at $p<0.05$ despite their size
(\cref{tab:sig}; $p$-values range 0.12--0.27). The reason is visible in
\cref{fig:rmse-box}: linear-model RMSE varies enormously across folds
(interquartile range $\sim$2500--6000 on \texttt{wd\_am}), which inflates
the test statistic's denominator. The size of the gap is real; the test
simply lacks power against this much fold-level variance with five
folds. We report the raw difference alongside the $p$-value for this
reason.

\paragraph{Ridge barely improves over OLS.}
Ridge reduces RMSE by only $2$--$4\%$ across targets, confirming that
multicollinearity is not the bottleneck for linear fit. The 13 features
come from different source tables and use different radii, so they are
reasonably de-correlated by design. Regularisation is not what linear
models need here --- more flexible functional form is.

\paragraph{Gradient boosting vs.\ random forest.}
GradBoost's edge over random forest is small in magnitude (97--247
RMSE) but more consistent across folds and reaches significance on
\texttt{wd\_am} ($p = 0.003$) and \texttt{h\_avg} ($p = 0.027$). For
\texttt{wd\_pm} and \texttt{wd\_midday} the difference does not reach
significance. The practical implication is that choice of tree
ensemble matters less than the linear-vs-tree decision: either
tree-based option closes nearly all of the gap against the linear
baseline.

\begin{table}[t]
  \centering
  \small
  \caption{Paired $t$-test comparing gradient boosting (best per target)
    against each alternative. $\Delta$RMSE is the excess RMSE of the
    alternative. Bold: significant at $p<0.05$.}
  \label{tab:sig}
  \begin{tabular}{@{}llcc@{}}
    \toprule
    Target              & Comparison     & $\Delta$RMSE & $p$-value        \\
    \midrule
    \texttt{wd\_am}     & GB vs.\ Linear & 1725.4       & 0.258            \\
                        & GB vs.\ Ridge  & 1600.9       & 0.267            \\
                        & GB vs.\ RF     & 153.8        & $\mathbf{0.003}$ \\
    \midrule
    \texttt{wd\_pm}     & GB vs.\ Linear & 1614.7       & 0.259            \\
                        & GB vs.\ Ridge  & 1498.1       & 0.267            \\
                        & GB vs.\ RF     & 247.1        & 0.104            \\
    \midrule
    \texttt{wd\_midday} & GB vs.\ Linear & 2338.5       & 0.123            \\
                        & GB vs.\ Ridge  & 2182.7       & 0.131            \\
                        & GB vs.\ RF     & 147.8        & 0.097            \\
    \midrule
    \texttt{h\_avg}     & GB vs.\ Linear & 695.5        & 0.160            \\
                        & GB vs.\ Ridge  & 644.7        & 0.171            \\
                        & GB vs.\ RF     & 96.8         & $\mathbf{0.027}$ \\
    \bottomrule
  \end{tabular}
\end{table}

\begin{figure}[t]
  \centering
  \includegraphics[width=0.95\columnwidth]{../reports/rmse_box}
  \caption{Per-fold test RMSE by model, one panel per target. Tree-based
    models (random forest, gradient boosting) are consistently far
    below linear models, but linear-model RMSE also varies dramatically
    across folds, which explains why the large gap does not reach
    statistical significance in a paired test.}
  \label{fig:rmse-box}
\end{figure}

\section{Stop Typology via Clustering}
\label{sec:clusters}

To complement the supervised analysis, we cluster stops on their
four-target usage profile (each stop's target vector normalised to sum
to 1) and ask whether natural groupings correspond to interpretable
transit archetypes. Normalising to profiles separates temporal
\emph{pattern} from overall \emph{busyness} --- an unnormalised clustering
on this data simply ranks stops by total volume, which is not what we
want. We fit K-means with $K = 4$, motivated by the regimes visible in
the hourly data (weekday morning peak, weekday evening peak, midday
off-peak, weekend) plus a natural split between morning-origin and
evening-destination stops.

\paragraph{Four archetypes emerge.}
\cref{tab:clusters} labels the resulting clusters.

\begin{table}[t]
  \centering
  \small
  \caption{K-means cluster archetypes ($K = 4$), fit on per-stop
    normalised target profiles.}
  \label{tab:clusters}
  \begin{tabular}{@{}cll@{}}
    \toprule
    $k$ & Profile                            & Interpretation        \\
    \midrule
    0   & AM $>$ PM $>$ midday $>$ weekend   & General residential   \\
    1   & Dominant AM, near-zero weekend     & Residential origin    \\
    2   & Dominant PM, near-zero weekend     & Workplace destination \\
    3   & Even AM / PM / midday, low weekend & Mixed-use             \\
    \bottomrule
  \end{tabular}
\end{table}

\paragraph{Geographic coherence.}
The clustering operates on temporal profiles only --- coordinates are
not used --- yet the archetypes lay out coherently on the map
(\cref{fig:cluster-geo}). Morning-origin and evening-destination stops
(clusters 1 and 2) concentrate in the western region around Jurong, the
part of Singapore with the clearest residential--industrial commute
structure: residents tap in during the morning rush at dormitory-style
estates and tap out in the evening at workplace clusters. Balanced and
mixed-use stops (clusters 0 and 3) dominate elsewhere, including the
core and the northeastern residential belt.

No cluster is dominated by weekend ridership; weekend volume is
uniformly modest across all archetypes. Temporal variation in this
network is primarily a weekday phenomenon, which also frames the
supervised result that weekend prediction is the hardest of the four
targets.

\begin{figure}[t]
  \centering
  \includegraphics[width=0.95\columnwidth]{../reports/cluster_geo}
  \caption{Singapore bus stops coloured by K-means cluster ($K = 4$),
    fit on normalised target profiles. Morning-origin (red, cluster 1)
    and evening-destination (pink, cluster 2) stops concentrate in the
    western Jurong region; balanced (blue, 0) and mixed-use (cyan, 3)
    archetypes dominate elsewhere.}
  \label{fig:cluster-geo}
\end{figure}

\section{Limitations}
\label{sec:limitations}

Several caveats qualify the findings.

\emph{Temporal scope.} LTSG is a snapshot from one period; ridership
patterns shift with urban development, COVID-era work-from-home
adjustments, and new MRT lines. Our features reflect current
infrastructure from LTA; a small mismatch between the two periods is
unavoidable.

\emph{Spatial aggregation choices.} The 500\,m / 1\,km radii are
defensible but not optimised.

\emph{Feature coverage.} We do not use POI quality signals (ratings, price
level) or bus-timing columns (first/last bus, which could yield a
service-span feature). Including these is a natural extension.

\emph{Independence assumption.} Our evaluation protocol treats stops as
i.i.d.\ samples, but stops close together share the same local
environment and likely have correlated volumes.

\emph{Four separate models.} Fitting each target independently ignores
the clear correlation between, e.g., \texttt{wd\_am} and \texttt{wd\_pm}.
A multi-output model could share strength and yield tighter estimates,
particularly for the smaller weekend signal.

\section{Conclusion}
\label{sec:conclusion}

\section*{Impact Statement}

This work uses only public, non-personal data (aggregated passenger counts,
infrastructure locations, and commercial amenity metadata). The intended
use is operational planning, where accurate demand estimation can improve
service efficiency and equity of access; failure modes (e.g.\ under-serving
stops with systematically underestimated volume) are worth checking
against ground truth before deployment. No individual-level data is
involved.

\bibliography{bib}
\bibliographystyle{icml2025}

\appendix
\onecolumn
\section{Appendix: Reproducibility}

Code and setup instructions are in the project repository. The pipeline is
three notebooks: \texttt{01\_eda.py} (data exploration and target
design), \texttt{02\_features.py} (feature engineering and join), and
\texttt{03\_model.py} (experiments, hyperparameter selection, and
plotting). All random seeds are fixed.

\end{document}
